\appendix
\section{Declaración detallada sobre el uso de inteligencia artificial}\label{anx:ia}
Se empleó un asistente de inteligencia artificial para (i) reorganizar y redactar el manuscrito, (ii) escribir y depurar el código Python de los cuadernos y de las figuras, (iii) analizar los archivos de resultados HDF de HEC-RAS (reparto, balance de volumen, resistencia efectiva) y (iv) preparar el generador de geometría 2D. El asistente no ejecutó las simulaciones HEC-RAS, que corrieron los integrantes, ni obtuvo observaciones de campo. Las fuentes y las citas incluidas deben ser verificadas por el grupo contra la fuente original, en especial \citet{gumgum2025} y \citet{hsu2002}.

\section{Estructura del repositorio y reproducción}
\begin{itemize}
\item \path{Art_MIH-626.tex}: archivo principal (compilar con pdflatex dos veces).
\item \path{water_like.sty}: estilo emulado de Water; \path{secciones/}, \path{referencias/}, \path{figuras/}.
\item \path{datos/}: \path{datos_entrada.json}, \path{resultados_2d_v3.csv}, \path{mallas_EM.csv}, \path{mallas_ELM.csv}, \path{elasticidades_fase4.csv}, \path{gci_calados_onda_difusiva.csv}.
\item \path{codigo/figuras_final.py}: regenera las figuras de reparto y malla; \path{codigo/figuras_articulo.py}: figuras de sensibilidad.
\end{itemize}

\section{Convergencia de calados máximos con onda difusiva (estudio preliminar)}\label{anx:gci}
Por sugerencia del docente se incluye el estudio de convergencia de malla que el grupo hizo con las ecuaciones de onda difusiva, anterior al paso a SWE-EM. La variable es el calado máximo en secciones de control de cada ramal, con celdas de 0.50, 0.25 y 0.125 m (razón de refinamiento 2) y el índice de convergencia de malla (GCI); la Tabla~\ref{tab:gci} resume los resultados. El umbral de 5\pc{} es una convención habitual del GCI y no un límite normativo.

\begin{table}[H]
\centering\small
\caption{Calados máximos [m] y GCI con onda difusiva (celdas de 0.50, 0.25 y 0.125 m).}\label{tab:gci}
\begin{tabular}{rcccccc}
\toprule
$\theta$ y ramal & $y$ (0.50) & $y$ (0.25) & $y$ (0.125) & $p$ & GCI 0.25 [\%] & GCI 0.125 [\%]\\\midrule
30°, R1 & 1.413 & 1.350 & 1.325 & 1.33 & 3.84 & 1.55\\
30°, R2 & 1.411 & 1.349 & 1.330 & 1.71 & 2.54 & 0.79\\
45°, R1$^{a}$ & 1.520 & 1.319 & 1.328 & --- & --- & ---\\
45°, R2$^{b}$ & 1.524 & 1.332 & 1.330 & --- & --- & ---\\
60°, R1 & 1.498 & 1.343 & 1.325 & 3.11 & 1.90 & 0.22\\
60°, R2 & 1.519 & 1.360 & 1.330 & 2.41 & 3.40 & 0.66\\
90°, R1 & 1.402 & 1.264 & 1.232 & 2.11 & 4.12 & 0.98\\
90°, R2 & 1.384 & 1.258 & 1.234 & 2.39 & 2.95 & 0.57\\\midrule
120°, R1 & 1.403 & 1.330 & 1.283 & 0.64 & 12.40 & 8.28\\
120°, R2 & 1.346 & 1.260 & 1.226 & 1.34 & 5.58 & 2.27\\\bottomrule
\end{tabular}

\smallskip\footnotesize $^{a}$ Convergencia no monótona (el calado sube al refinar): el GCI no es aplicable; la hoja del grupo informó $p=4.48$ tomando valores absolutos. $^{b}$ Diferencia de 2 mm entre las dos mallas finas: $p$ no es fiable (la hoja informó 6.58). Los valores de 120° se presentan aparte porque sus resultados no son consistentes con los demás.
\end{table}

Con celda de 0.25 m, el GCI de los calados es de a lo sumo 4.1\pc{} entre 30° y 90° en las secciones donde la convergencia es monótona, y la celda de 0.50 m queda descartada porque los calados difieren 14 a 20 cm de los de 0.25 m. A 120° el GCI supera el umbral (hasta 12.4\pc), lo que el grupo atribuyó a los giros cerrados del flujo y motivó el uso de SWE-EM. Estos resultados son con onda difusiva y de calados máximos: \textbf{no demuestran independencia de malla del reparto $\varphi$} con SWE-EM, que sí cambia con la celda (Tabla~\ref{tab:mallas}).

\section{Cumplimiento de los requisitos del trabajo}
\begin{table}[H]
\centering\small
\caption{Estado de los requisitos del Trabajo 01.}\label{tab:req}
\begin{tabular}{>{\raggedright\arraybackslash}p{0.40\linewidth}>{\raggedright\arraybackslash}p{0.16\linewidth}>{\raggedright\arraybackslash}p{0.34\linewidth}}
\toprule
Requisito & Estado & Dónde / observación\\\midrule
Sistema no lineal resuelto por Newton--Raphson & Cumplido & Secciones~\ref{sec:f1}--\ref{sec:f2}; \path{codigo/newton_bifurcacion.py}\\
Tres pruebas de verificación & Cumplido, una no superada & Tabla~\ref{tab:verif}; el benchmark de Chaudhry no se reproduce\\
Comparación con HEC-RAS 2D & Cumplido & Tabla~\ref{tab:reparto}; cinco ángulos con SWE-EM y celda de 0.25 m\\
Tres mallas y balance de masa & Balance sí; mallas parcial & Tablas~\ref{tab:cc} y~\ref{tab:mallas}; tres mallas solo en 30° y 120° (SWE-EM) y en los cinco ángulos (SWE-ELM, exploratorio); GCI de calados con onda difusiva en el Anexo~\ref{anx:gci}\\
Sensibilidad (al menos cinco parámetros) & Cumplido & 13 parámetros; Figura~\ref{fig:elast}; barrido de $K$ y $n$ en la Tabla~\ref{tab:barrido}\\
Discusión de discrepancias & Cumplido & Sección~\ref{sec:discusion}\\
Datos en archivo separado, versiones impresas, semilla fija & Cumplido & \path{datos/}; los \emph{scripts} imprimen versiones; semilla 20261103\\
Al menos 15 referencias, 5 arbitradas, formato APA & Cumplido, por verificar & Referencias; verificar los datos de \citet{gumgum2025} y \citet{hsu2002}\\
Declaración de uso de IA & Cumplido & Anexo~\ref{anx:ia}\\
Repositorio & Por actualizar & \repo; subir esta carpeta y los anexos de HEC-RAS\\
Informe de 15 a 25 páginas & Verificar & Depende de la compilación final\\\bottomrule
\end{tabular}
\end{table}

\section{Inventario de archivos de HEC-RAS utilizados}
\begin{itemize}
\item Paquetes de geometría v3 por ángulo (\texttt{theta30}, \dots, \texttt{theta120}): terreno, área 2D, líneas de quiebre, líneas de borde, regiones de Manning, hidrograma de entrada y niveles de salida.
\item Corridas SWE-EM: 30° y 120° con mallas de 0.50, 0.25 y 0.125 m y refinamientos del nudo; 45°, 60° y 90° con malla de 0.25 m.
\item Ramales rectos aislados (ramal 1 y ramal 2) con mallas de 0.50, 0.25 y 0.125 m y los dos planes con $n$ corregido.
\item Conjunto exploratorio SWE-ELM: resultados CSV por ángulo y malla (\path{ProcesamientoResultados/parac0.5}, \path{parac0.25}, \path{parac0.125}).
\end{itemize}

\section{Pasos para reproducir los resultados}
\begin{enumerate}
\item Ejecutar \texttt{python codigo/newton\_bifurcacion.py}: imprime versiones, historial de Newton, $\varphi_{\mathrm{II}}$ y $\varphi_{\mathrm I}$.
\item Ejecutar \texttt{python codigo/figuras\_final.py} y \texttt{python codigo/figuras\_articulo.py}: regeneran las figuras desde \path{datos/}.
\item Compilar \path{Art_MIH-626.tex} con pdflatex dos veces (no se requiere bibtex ni biber).
\item Los cuadernos de la Fase IV (sesgo y sensibilidad) y del generador 2D se conservan en el repositorio general del trabajo; sus salidas son la fuente de \path{datos/resultados_fase4.csv}.
\end{enumerate}
